% Options for packages loaded elsewhere
% Options for packages loaded elsewhere
\PassOptionsToPackage{unicode}{hyperref}
\PassOptionsToPackage{hyphens}{url}
\PassOptionsToPackage{dvipsnames,svgnames,x11names}{xcolor}
%
\documentclass[
  letterpaper,
  DIV=11,
  numbers=noendperiod]{scrartcl}
\usepackage{xcolor}
\usepackage{amsmath,amssymb}
\setcounter{secnumdepth}{-\maxdimen} % remove section numbering
\usepackage{iftex}
\ifPDFTeX
  \usepackage[T1]{fontenc}
  \usepackage[utf8]{inputenc}
  \usepackage{textcomp} % provide euro and other symbols
\else % if luatex or xetex
  \usepackage{unicode-math} % this also loads fontspec
  \defaultfontfeatures{Scale=MatchLowercase}
  \defaultfontfeatures[\rmfamily]{Ligatures=TeX,Scale=1}
\fi
\usepackage{lmodern}
\ifPDFTeX\else
  % xetex/luatex font selection
\fi
% Use upquote if available, for straight quotes in verbatim environments
\IfFileExists{upquote.sty}{\usepackage{upquote}}{}
\IfFileExists{microtype.sty}{% use microtype if available
  \usepackage[]{microtype}
  \UseMicrotypeSet[protrusion]{basicmath} % disable protrusion for tt fonts
}{}
\makeatletter
\@ifundefined{KOMAClassName}{% if non-KOMA class
  \IfFileExists{parskip.sty}{%
    \usepackage{parskip}
  }{% else
    \setlength{\parindent}{0pt}
    \setlength{\parskip}{6pt plus 2pt minus 1pt}}
}{% if KOMA class
  \KOMAoptions{parskip=half}}
\makeatother
% Make \paragraph and \subparagraph free-standing
\makeatletter
\ifx\paragraph\undefined\else
  \let\oldparagraph\paragraph
  \renewcommand{\paragraph}{
    \@ifstar
      \xxxParagraphStar
      \xxxParagraphNoStar
  }
  \newcommand{\xxxParagraphStar}[1]{\oldparagraph*{#1}\mbox{}}
  \newcommand{\xxxParagraphNoStar}[1]{\oldparagraph{#1}\mbox{}}
\fi
\ifx\subparagraph\undefined\else
  \let\oldsubparagraph\subparagraph
  \renewcommand{\subparagraph}{
    \@ifstar
      \xxxSubParagraphStar
      \xxxSubParagraphNoStar
  }
  \newcommand{\xxxSubParagraphStar}[1]{\oldsubparagraph*{#1}\mbox{}}
  \newcommand{\xxxSubParagraphNoStar}[1]{\oldsubparagraph{#1}\mbox{}}
\fi
\makeatother

\usepackage{color}
\usepackage{fancyvrb}
\newcommand{\VerbBar}{|}
\newcommand{\VERB}{\Verb[commandchars=\\\{\}]}
\DefineVerbatimEnvironment{Highlighting}{Verbatim}{commandchars=\\\{\}}
% Add ',fontsize=\small' for more characters per line
\usepackage{framed}
\definecolor{shadecolor}{RGB}{241,243,245}
\newenvironment{Shaded}{\begin{snugshade}}{\end{snugshade}}
\newcommand{\AlertTok}[1]{\textcolor[rgb]{0.68,0.00,0.00}{#1}}
\newcommand{\AnnotationTok}[1]{\textcolor[rgb]{0.37,0.37,0.37}{#1}}
\newcommand{\AttributeTok}[1]{\textcolor[rgb]{0.40,0.45,0.13}{#1}}
\newcommand{\BaseNTok}[1]{\textcolor[rgb]{0.68,0.00,0.00}{#1}}
\newcommand{\BuiltInTok}[1]{\textcolor[rgb]{0.00,0.23,0.31}{#1}}
\newcommand{\CharTok}[1]{\textcolor[rgb]{0.13,0.47,0.30}{#1}}
\newcommand{\CommentTok}[1]{\textcolor[rgb]{0.37,0.37,0.37}{#1}}
\newcommand{\CommentVarTok}[1]{\textcolor[rgb]{0.37,0.37,0.37}{\textit{#1}}}
\newcommand{\ConstantTok}[1]{\textcolor[rgb]{0.56,0.35,0.01}{#1}}
\newcommand{\ControlFlowTok}[1]{\textcolor[rgb]{0.00,0.23,0.31}{\textbf{#1}}}
\newcommand{\DataTypeTok}[1]{\textcolor[rgb]{0.68,0.00,0.00}{#1}}
\newcommand{\DecValTok}[1]{\textcolor[rgb]{0.68,0.00,0.00}{#1}}
\newcommand{\DocumentationTok}[1]{\textcolor[rgb]{0.37,0.37,0.37}{\textit{#1}}}
\newcommand{\ErrorTok}[1]{\textcolor[rgb]{0.68,0.00,0.00}{#1}}
\newcommand{\ExtensionTok}[1]{\textcolor[rgb]{0.00,0.23,0.31}{#1}}
\newcommand{\FloatTok}[1]{\textcolor[rgb]{0.68,0.00,0.00}{#1}}
\newcommand{\FunctionTok}[1]{\textcolor[rgb]{0.28,0.35,0.67}{#1}}
\newcommand{\ImportTok}[1]{\textcolor[rgb]{0.00,0.46,0.62}{#1}}
\newcommand{\InformationTok}[1]{\textcolor[rgb]{0.37,0.37,0.37}{#1}}
\newcommand{\KeywordTok}[1]{\textcolor[rgb]{0.00,0.23,0.31}{\textbf{#1}}}
\newcommand{\NormalTok}[1]{\textcolor[rgb]{0.00,0.23,0.31}{#1}}
\newcommand{\OperatorTok}[1]{\textcolor[rgb]{0.37,0.37,0.37}{#1}}
\newcommand{\OtherTok}[1]{\textcolor[rgb]{0.00,0.23,0.31}{#1}}
\newcommand{\PreprocessorTok}[1]{\textcolor[rgb]{0.68,0.00,0.00}{#1}}
\newcommand{\RegionMarkerTok}[1]{\textcolor[rgb]{0.00,0.23,0.31}{#1}}
\newcommand{\SpecialCharTok}[1]{\textcolor[rgb]{0.37,0.37,0.37}{#1}}
\newcommand{\SpecialStringTok}[1]{\textcolor[rgb]{0.13,0.47,0.30}{#1}}
\newcommand{\StringTok}[1]{\textcolor[rgb]{0.13,0.47,0.30}{#1}}
\newcommand{\VariableTok}[1]{\textcolor[rgb]{0.07,0.07,0.07}{#1}}
\newcommand{\VerbatimStringTok}[1]{\textcolor[rgb]{0.13,0.47,0.30}{#1}}
\newcommand{\WarningTok}[1]{\textcolor[rgb]{0.37,0.37,0.37}{\textit{#1}}}

\usepackage{longtable,booktabs,array}
\usepackage{calc} % for calculating minipage widths
% Correct order of tables after \paragraph or \subparagraph
\usepackage{etoolbox}
\makeatletter
\patchcmd\longtable{\par}{\if@noskipsec\mbox{}\fi\par}{}{}
\makeatother
% Allow footnotes in longtable head/foot
\IfFileExists{footnotehyper.sty}{\usepackage{footnotehyper}}{\usepackage{footnote}}
\makesavenoteenv{longtable}
\usepackage{graphicx}
\makeatletter
\newsavebox\pandoc@box
\newcommand*\pandocbounded[1]{% scales image to fit in text height/width
  \sbox\pandoc@box{#1}%
  \Gscale@div\@tempa{\textheight}{\dimexpr\ht\pandoc@box+\dp\pandoc@box\relax}%
  \Gscale@div\@tempb{\linewidth}{\wd\pandoc@box}%
  \ifdim\@tempb\p@<\@tempa\p@\let\@tempa\@tempb\fi% select the smaller of both
  \ifdim\@tempa\p@<\p@\scalebox{\@tempa}{\usebox\pandoc@box}%
  \else\usebox{\pandoc@box}%
  \fi%
}
% Set default figure placement to htbp
\def\fps@figure{htbp}
\makeatother





\setlength{\emergencystretch}{3em} % prevent overfull lines

\providecommand{\tightlist}{%
  \setlength{\itemsep}{0pt}\setlength{\parskip}{0pt}}



 


\KOMAoption{captions}{tableheading}
\makeatletter
\@ifpackageloaded{caption}{}{\usepackage{caption}}
\AtBeginDocument{%
\ifdefined\contentsname
  \renewcommand*\contentsname{Table of contents}
\else
  \newcommand\contentsname{Table of contents}
\fi
\ifdefined\listfigurename
  \renewcommand*\listfigurename{List of Figures}
\else
  \newcommand\listfigurename{List of Figures}
\fi
\ifdefined\listtablename
  \renewcommand*\listtablename{List of Tables}
\else
  \newcommand\listtablename{List of Tables}
\fi
\ifdefined\figurename
  \renewcommand*\figurename{Figure}
\else
  \newcommand\figurename{Figure}
\fi
\ifdefined\tablename
  \renewcommand*\tablename{Table}
\else
  \newcommand\tablename{Table}
\fi
}
\@ifpackageloaded{float}{}{\usepackage{float}}
\floatstyle{ruled}
\@ifundefined{c@chapter}{\newfloat{codelisting}{h}{lop}}{\newfloat{codelisting}{h}{lop}[chapter]}
\floatname{codelisting}{Listing}
\newcommand*\listoflistings{\listof{codelisting}{List of Listings}}
\makeatother
\makeatletter
\makeatother
\makeatletter
\@ifpackageloaded{caption}{}{\usepackage{caption}}
\@ifpackageloaded{subcaption}{}{\usepackage{subcaption}}
\makeatother
\usepackage{bookmark}
\IfFileExists{xurl.sty}{\usepackage{xurl}}{} % add URL line breaks if available
\urlstyle{same}
\hypersetup{
  pdftitle={Tarea 1},
  colorlinks=true,
  linkcolor={blue},
  filecolor={Maroon},
  citecolor={Blue},
  urlcolor={Blue},
  pdfcreator={LaTeX via pandoc}}


\title{Tarea 1}
\author{}
\date{}
\begin{document}
\maketitle


\subsection{Introducción a la Ciencia de Datos
2025}\label{introducciuxf3n-a-la-ciencia-de-datos-2025}

\subsubsection{Ignacio Campón \& Mauro
Loprete}\label{ignacio-campuxf3n-mauro-loprete}

\begin{Shaded}
\begin{Highlighting}[]
\ImportTok{from}\NormalTok{ time }\ImportTok{import}\NormalTok{ time}
\ImportTok{from}\NormalTok{ pathlib }\ImportTok{import}\NormalTok{ Path}

\ImportTok{import}\NormalTok{ pandas }\ImportTok{as}\NormalTok{ pd}
\ImportTok{import}\NormalTok{ matplotlib.pyplot }\ImportTok{as}\NormalTok{ plt}
\ImportTok{import}\NormalTok{ seaborn }\ImportTok{as}\NormalTok{ sns}
\ImportTok{import}\NormalTok{ networkx }\ImportTok{as}\NormalTok{ nx}
\ImportTok{import}\NormalTok{ numpy }\ImportTok{as}\NormalTok{ np}
\ImportTok{from}\NormalTok{ tabulate }\ImportTok{import}\NormalTok{ tabulate}
\ImportTok{import}\NormalTok{ re}

\ImportTok{from}\NormalTok{ collections }\ImportTok{import}\NormalTok{ Counter}
\ImportTok{from}\NormalTok{ nltk.corpus }\ImportTok{import}\NormalTok{ stopwords}

\ImportTok{import}\NormalTok{ nltk}

\NormalTok{sns.set\_theme(style}\OperatorTok{=}\StringTok{"whitegrid"}\NormalTok{)}
\end{Highlighting}
\end{Shaded}

\subsection{Parte 1: Cargado y Limpieza de
datos}\label{parte-1-cargado-y-limpieza-de-datos}

\begin{Shaded}
\begin{Highlighting}[]
\NormalTok{df\_speeches }\OperatorTok{=}\NormalTok{ pd.read\_csv(}\StringTok{\textquotesingle{}../data/us\_2020\_election\_speeches.csv\textquotesingle{}}\NormalTok{)}

\NormalTok{df\_speeches}
\end{Highlighting}
\end{Shaded}

\begin{longtable}[]{@{}lllllll@{}}
\toprule\noalign{}
& speaker & title & text & date & location & type \\
\midrule\noalign{}
\endhead
\bottomrule\noalign{}
\endlastfoot
0 & David Perdue & Georgia Sen. David Perdue Speech Transcript at... &
David Perdue: (00:01)\textbackslash nHow great is it to be b... & Oct
16, 2020 & Macon, Georgia & Campaign Speech \\
1 & Joe Biden & Joe Biden Southfield, MI Speech on Health Care... & Joe
Biden: (00:00)\textbackslash nHello, Michigan. Hi, how a... & Oct 16,
2020 & Southfield ,Michigan & Campaign Speech \\
2 & Donald Trump & Donald Trump Speech Transcript `Protecting Ame... &
President Trump: (00:30)\textbackslash nThank you. What a ni... & Oct
16, 2020 & Fort Myers, Florida & Campaign Speech \\
3 & Joe Biden & Joe Biden ABC Town Hall Transcript October 15 & George
Stephanopoulos: (00:41)\textbackslash nHey, and welco... & Oct 15, 2020
& ABC & Town Hall \\
4 & Donald Trump & Donald Trump NBC Town Hall Transcript October 15 &
Savannah Guthrie: (03:50)\textbackslash nIt's nothing but no... & Oct
15, 2020 & NBC & Town Hall \\
... & ... & ... & ... & ... & ... & ... \\
264 & Bernie Sanders & Bernie Sanders Speech Transcript: Sanders Spea...
& Bernie Sanders: (00:00)\textbackslash nJust want to take thi... & Feb
6, 2020 & Iowa & Campaign Speech \\
265 & Democratic Candidates & Transcript: Speeches at the Iowa Caucuses
-- Be... & Bernie Sanders: (00:08)\textbackslash nThank you. Thank
you.... & Feb 4, 2020 & Iowa & Campaign Speech \\
266 & Donald Trump & Donal Trump Iowa Rally Transcript: Trump Holds... &
Donald Trump: (00:24)\textbackslash nI worked so hard for th... & Jan
30, 2020 & Des Moines, Iowa & Campaign Speech \\
267 & Donald Trump & Donald Trump New Jersey Rally Speech Transcrip... &
Donald Trump: (01:22)\textbackslash nThank you. Thank you. I... & Jan
28, 2020 & Wildwood, New Jersey & Campaign Speech \\
268 & Democratic Candidates & January Iowa Democratic Debate Transcript
& Wolf Blitzer: (00:00)\textbackslash nAll right, so let's beg... & Jan
15, 2020 & Des Moines, Iowa & Debate \\
\end{longtable}

\section{A) Datos faltantes}\label{a-datos-faltantes}

La base de datos se compone de un conjunto de 269 observaciones y 6
variables que contienen discursos de campaña, discursos de ayuntamiento,
entrevistas, debates y más. Los discursos estan ingresados a la base por
orador, título, contenido hablado, fecha, localidad y tipo de discuro.
En particular se analizaran los discursos de campaña política, y se hará
enfásis en los 5 prinipales oradores que seán seleccionados por mayor
frecuencia de discursos.

Para seleccionar esta categoría de discursos, primero se realiza un
análisis de datos faltantes en la base. Para ello primero verificamos
los tipos de variables con los que se cuentan los cuáles son 3 tipos.
Variables categóricas; \emph{speaker}, describe el nombre del orador,
\emph{type} es el tipo de discurso y \emph{location} la ciudad y el
estado en el que se hizo. Texto; \emph{title} título del discruso y
\emph{text} contenido del mismo. Fecha: \emph{date} contiene la fecha en
formato año/dia/mes . Los datos faltantes por variable se describen en
la tabla **\_\_**:

\begin{longtable}[]{@{}ll@{}}
\toprule\noalign{}
Variable & Cantidad de Nulos \\
\midrule\noalign{}
\endhead
\bottomrule\noalign{}
\endlastfoot
speaker & 3 \\
title & 0 \\
text & 0 \\
date & 0 \\
location & 18 \\
type & 21 \\
\end{longtable}

Para verificar la calidad de los datos, realizamos una visualización
manual en busca de datos incompletos, datos inconsistentes, datos con
ruido que probablemente sean provenientes de una mala producción y/o
procesamiento de los mismos. La busqueda no arroja mala calidad de los
mismos, resultando haber una buena calidad de los mismos.

\begin{Shaded}
\begin{Highlighting}[]
\CommentTok{\# visualizamos los datos faltantes por variable}
\NormalTok{n }\OperatorTok{=} \BuiltInTok{len}\NormalTok{(df\_speeches)}
\NormalTok{count }\OperatorTok{=}\NormalTok{ df\_speeches.count()}
\NormalTok{res }\OperatorTok{=}\NormalTok{ n}\OperatorTok{{-}}\NormalTok{count}

\CommentTok{\# creamos tabla con los datos faltantes}
\NormalTok{n\_count\_df }\OperatorTok{=}\NormalTok{ res.reset\_index()}
\NormalTok{n\_count\_df.columns }\OperatorTok{=}\NormalTok{ [}\StringTok{\textquotesingle{}Variable\textquotesingle{}}\NormalTok{, }\StringTok{\textquotesingle{}Cantidad de Nulos\textquotesingle{}}\NormalTok{]}

\CommentTok{\#print(tabulate(n\_count\_df, headers=\textquotesingle{}keys\textquotesingle{}, tablefmt=\textquotesingle{}github\textquotesingle{}, showindex=False))}
\end{Highlighting}
\end{Shaded}

\begin{Shaded}
\begin{Highlighting}[]
\CommentTok{\#df\_speeches.dtypes}
\CommentTok{\#df\_speeches.describe()}
\end{Highlighting}
\end{Shaded}

\begin{Shaded}
\begin{Highlighting}[]
\NormalTok{df\_speeches[}\StringTok{"date"}\NormalTok{] }\OperatorTok{=}\NormalTok{ pd.to\_datetime(df\_speeches[}\StringTok{"date"}\NormalTok{])}

\KeywordTok{def}\NormalTok{ get\_var\_category(series):}
\NormalTok{    unique\_count }\OperatorTok{=}\NormalTok{ series.nunique(dropna}\OperatorTok{=}\VariableTok{False}\NormalTok{)}
\NormalTok{    total\_count }\OperatorTok{=} \BuiltInTok{len}\NormalTok{(series)}
    \ControlFlowTok{if}\NormalTok{ pd.api.types.is\_numeric\_dtype(series):}
        \ControlFlowTok{return} \StringTok{\textquotesingle{}Numerical\textquotesingle{}}
    \ControlFlowTok{elif}\NormalTok{ pd.api.types.is\_datetime64\_dtype(series):}
        \ControlFlowTok{return} \StringTok{\textquotesingle{}Date\textquotesingle{}}
    \ControlFlowTok{elif}\NormalTok{ unique\_count}\OperatorTok{==}\NormalTok{total\_count:}
        \ControlFlowTok{return} \StringTok{\textquotesingle{}Text (Unique)\textquotesingle{}}
    \ControlFlowTok{else}\NormalTok{:}
        \ControlFlowTok{return} \StringTok{\textquotesingle{}Categorical\textquotesingle{}}

\KeywordTok{def}\NormalTok{ print\_categories(df):}
    \ControlFlowTok{for}\NormalTok{ column\_name }\KeywordTok{in}\NormalTok{ df.columns:}
        \BuiltInTok{print}\NormalTok{(column\_name, }\StringTok{": "}\NormalTok{, get\_var\_category(df[column\_name]))}

\CommentTok{\#print\_categories(df\_speeches)}
\end{Highlighting}
\end{Shaded}

\begin{Shaded}
\begin{Highlighting}[]
\CommentTok{\# how many types and obs are? creamos tabla con los tipos de discurso}
\NormalTok{type\_counts }\OperatorTok{=}\NormalTok{ df\_speeches[}\StringTok{"type"}\NormalTok{].value\_counts().reset\_index()}
\NormalTok{type\_counts.columns }\OperatorTok{=}\NormalTok{ [}\StringTok{"Tipo de Discurso"}\NormalTok{, }\StringTok{"Cantidad"}\NormalTok{]}

\CommentTok{\#print(tabulate(type\_counts, headers="keys", tablefmt="github", showindex=False))}
\end{Highlighting}
\end{Shaded}

La cantidad de observaciones por \emph{type} de discurso se resumen en
la tabla **\_\_**. Se puede observar que 180 de las 269 observaciones
corresponden a discursos de campaña los cuáles serán nuestra muestra de
análisis.

\begin{longtable}[]{@{}ll@{}}
\toprule\noalign{}
Tipo de Discurso & Cantidad \\
\midrule\noalign{}
\endhead
\bottomrule\noalign{}
\endlastfoot
Campaign Speech & 180 \\
Town Hall & 18 \\
Interview & 14 \\
Debate & 9 \\
Endorsement & 8 \\
Statement & 8 \\
Roundtable & 8 \\
Press Conference & 2 \\
Voter Mobilization & 1 \\
\end{longtable}

A modo informativo, los discursos de campaña toman fechas que van desde
el 15 de Enero de 2020 al 16 de Octubre de 2020, y la mediana de los
mismos toma fecha en fines de Agosto, dando pie a que la mayoria de los
mismos toma parte en los últimos meses de campaña.

\begin{Shaded}
\begin{Highlighting}[]
\CommentTok{\#df\_speeches.describe()}
\end{Highlighting}
\end{Shaded}

\begin{Shaded}
\begin{Highlighting}[]
\CommentTok{\# filter campaign speeches}
\NormalTok{speeches }\OperatorTok{=}\NormalTok{ df\_speeches[df\_speeches[}\StringTok{"type"}\NormalTok{] }\OperatorTok{==} \StringTok{"Campaign Speech"}\NormalTok{]}

\CommentTok{\# group speeches by speaker}
\CommentTok{\#speeches["speaker"].value\_counts()}

\CommentTok{\# group speeches by speaker and date}
\CommentTok{\#speeches.groupby(["date", "speaker"]).value\_counts()}

\CommentTok{\# filter speeches by donald trump}
\CommentTok{\#speeches[speeches["speaker"] == "Donald Trump"]}

\CommentTok{\#print(speeches.describe())}

\NormalTok{top5\_speakers }\OperatorTok{=}\NormalTok{ speeches[}\StringTok{"speaker"}\NormalTok{].value\_counts().head(}\DecValTok{5}\NormalTok{).index}
\NormalTok{top\_speeches }\OperatorTok{=}\NormalTok{ speeches[speeches[}\StringTok{"speaker"}\NormalTok{].isin(top5\_speakers)]}
\end{Highlighting}
\end{Shaded}

\begin{Shaded}
\begin{Highlighting}[]
\CommentTok{\# verificar si hay valores nulos después de la conversión}
\ControlFlowTok{if}\NormalTok{ top\_speeches[}\StringTok{"date"}\NormalTok{].isnull().}\BuiltInTok{any}\NormalTok{():}
    \BuiltInTok{print}\NormalTok{(}\StringTok{"Advertencia: Hay valores nulos en la columna \textquotesingle{}date\textquotesingle{} después de la conversión."}\NormalTok{)}

\CommentTok{\# agrupamos por mes y por candidato}
\NormalTok{top\_speeches[}\StringTok{"month"}\NormalTok{] }\OperatorTok{=}\NormalTok{ top\_speeches[}\StringTok{"date"}\NormalTok{].dt.to\_period(}\StringTok{"M"}\NormalTok{)}
\NormalTok{top\_speeches\_grouped }\OperatorTok{=}\NormalTok{ top\_speeches.groupby([}\StringTok{"month"}\NormalTok{, }\StringTok{"speaker"}\NormalTok{]).size().reset\_index(name}\OperatorTok{=}\StringTok{"count"}\NormalTok{)}

\CommentTok{\# Convertir el período a una fecha para el gráfico}
\NormalTok{top\_speeches\_grouped[}\StringTok{"month\_plot"}\NormalTok{] }\OperatorTok{=}\NormalTok{ top\_speeches\_grouped[}\StringTok{"month"}\NormalTok{].dt.to\_timestamp()}

\CommentTok{\# agregamos 0 a los meses que no tienen discursos}
\NormalTok{top\_speeches[}\StringTok{"date"}\NormalTok{] }\OperatorTok{=}\NormalTok{ pd.to\_datetime(top\_speeches[}\StringTok{"date"}\NormalTok{], errors}\OperatorTok{=}\StringTok{"coerce"}\NormalTok{)}
\NormalTok{top\_speeches[}\StringTok{"month"}\NormalTok{] }\OperatorTok{=}\NormalTok{ top\_speeches[}\StringTok{"date"}\NormalTok{].dt.to\_period(}\StringTok{"M"}\NormalTok{)}
\NormalTok{grouped }\OperatorTok{=}\NormalTok{ top\_speeches.groupby([}\StringTok{"month"}\NormalTok{, }\StringTok{"speaker"}\NormalTok{]).size().reset\_index(name}\OperatorTok{=}\StringTok{"count"}\NormalTok{)}
\NormalTok{grouped[}\StringTok{"month\_plot"}\NormalTok{] }\OperatorTok{=}\NormalTok{ grouped[}\StringTok{"month"}\NormalTok{].dt.to\_timestamp()}

\NormalTok{all\_months }\OperatorTok{=}\NormalTok{ pd.period\_range(start}\OperatorTok{=}\NormalTok{grouped[}\StringTok{"month"}\NormalTok{].}\BuiltInTok{min}\NormalTok{(), end}\OperatorTok{=}\NormalTok{grouped[}\StringTok{"month"}\NormalTok{].}\BuiltInTok{max}\NormalTok{(), freq}\OperatorTok{=}\StringTok{"M"}\NormalTok{)}
\NormalTok{all\_speakers }\OperatorTok{=}\NormalTok{ top\_speeches[}\StringTok{"speaker"}\NormalTok{].unique()}

\NormalTok{full\_index }\OperatorTok{=}\NormalTok{ pd.MultiIndex.from\_product([all\_months, all\_speakers], names}\OperatorTok{=}\NormalTok{[}\StringTok{"month"}\NormalTok{, }\StringTok{"speaker"}\NormalTok{])}
\NormalTok{full\_df }\OperatorTok{=}\NormalTok{ pd.DataFrame(index}\OperatorTok{=}\NormalTok{full\_index).reset\_index()}
\NormalTok{full\_merged }\OperatorTok{=}\NormalTok{ pd.merge(full\_df, grouped, on}\OperatorTok{=}\NormalTok{[}\StringTok{"month"}\NormalTok{, }\StringTok{"speaker"}\NormalTok{], how}\OperatorTok{=}\StringTok{"left"}\NormalTok{)}
\NormalTok{full\_merged[}\StringTok{"count"}\NormalTok{] }\OperatorTok{=}\NormalTok{ full\_merged[}\StringTok{"count"}\NormalTok{].fillna(}\DecValTok{0}\NormalTok{).astype(}\BuiltInTok{int}\NormalTok{)}
\NormalTok{full\_merged[}\StringTok{"month\_plot"}\NormalTok{] }\OperatorTok{=}\NormalTok{ full\_merged[}\StringTok{"month"}\NormalTok{].dt.to\_timestamp()}

\NormalTok{custom\_palette }\OperatorTok{=}\NormalTok{ [}\StringTok{"\#E69F00"}\NormalTok{, }\StringTok{"\#56B4E9"}\NormalTok{, }\StringTok{"\#009E73"}\NormalTok{, }\StringTok{"\#F0E442"}\NormalTok{, }\StringTok{"\#0072B2"}\NormalTok{]}



\CommentTok{\# Agrupamos por candidato y calculamos el total}
\NormalTok{lollipop\_data }\OperatorTok{=}\NormalTok{ speeches.groupby(}\StringTok{"speaker"}\NormalTok{).size().reset\_index(name}\OperatorTok{=}\StringTok{"count"}\NormalTok{)}

\CommentTok{\# Ordenamos por cantidad de discursos}
\NormalTok{lollipop\_data }\OperatorTok{=}\NormalTok{ lollipop\_data.sort\_values(by}\OperatorTok{=}\StringTok{"count"}\NormalTok{, ascending}\OperatorTok{=}\VariableTok{False}\NormalTok{)}

\CommentTok{\# Calculamos el total y los acumulados}
\NormalTok{total\_count }\OperatorTok{=}\NormalTok{ lollipop\_data[}\StringTok{"count"}\NormalTok{].}\BuiltInTok{sum}\NormalTok{()}
\NormalTok{lollipop\_data[}\StringTok{"percentage"}\NormalTok{] }\OperatorTok{=}\NormalTok{ lollipop\_data[}\StringTok{"count"}\NormalTok{] }\OperatorTok{/}\NormalTok{ total\_count}
\NormalTok{lollipop\_data[}\StringTok{"cumulative\_percentage"}\NormalTok{] }\OperatorTok{=}\NormalTok{ lollipop\_data[}\StringTok{"percentage"}\NormalTok{].cumsum()}

\CommentTok{\# Agregamos la fila "Otros"}
\NormalTok{top5\_data }\OperatorTok{=}\NormalTok{ lollipop\_data.head(}\DecValTok{5}\NormalTok{)}
\NormalTok{otros\_count }\OperatorTok{=}\NormalTok{ total\_count }\OperatorTok{{-}}\NormalTok{ top5\_data[}\StringTok{"count"}\NormalTok{].}\BuiltInTok{sum}\NormalTok{()}
\NormalTok{otros\_row }\OperatorTok{=}\NormalTok{ pd.DataFrame([[}\StringTok{"Otros"}\NormalTok{, otros\_count, otros\_count }\OperatorTok{/}\NormalTok{ total\_count, }\FloatTok{1.0}\NormalTok{]], }
\NormalTok{                         columns}\OperatorTok{=}\NormalTok{[}\StringTok{"speaker"}\NormalTok{, }\StringTok{"count"}\NormalTok{, }\StringTok{"percentage"}\NormalTok{, }\StringTok{"cumulative\_percentage"}\NormalTok{])}
\NormalTok{lollipop\_data\_extended }\OperatorTok{=}\NormalTok{ pd.concat([top5\_data, otros\_row], ignore\_index}\OperatorTok{=}\VariableTok{True}\NormalTok{)}

\NormalTok{chocolate\_data }\OperatorTok{=}\NormalTok{ lollipop\_data\_extended.copy()}
\NormalTok{chocolate\_data }\OperatorTok{=}\NormalTok{ chocolate\_data.sort\_values(by}\OperatorTok{=}\StringTok{"count"}\NormalTok{, ascending}\OperatorTok{=}\VariableTok{False}\NormalTok{)  }

\NormalTok{otros\_count }\OperatorTok{=}\NormalTok{ speeches[}\StringTok{"speaker"}\NormalTok{].nunique() }\OperatorTok{{-}} \DecValTok{5} 

\NormalTok{chocolate\_data[}\StringTok{"cumulative\_percentage"}\NormalTok{] }\OperatorTok{=}\NormalTok{ chocolate\_data[}\StringTok{"count"}\NormalTok{].cumsum() }\OperatorTok{/}\NormalTok{ chocolate\_data[}\StringTok{"count"}\NormalTok{].}\BuiltInTok{sum}\NormalTok{()}


\NormalTok{plt.figure(figsize}\OperatorTok{=}\NormalTok{(}\DecValTok{12}\NormalTok{, }\DecValTok{8}\NormalTok{))}
\NormalTok{bars }\OperatorTok{=}\NormalTok{ plt.barh(}
\NormalTok{    chocolate\_data[}\StringTok{"speaker"}\NormalTok{],}
\NormalTok{    chocolate\_data[}\StringTok{"count"}\NormalTok{],}
\NormalTok{    color}\OperatorTok{=}\NormalTok{custom\_palette, }
\NormalTok{    edgecolor}\OperatorTok{=}\StringTok{"black"}
\NormalTok{)}

\NormalTok{threshold }\OperatorTok{=} \FloatTok{0.654}

\NormalTok{cumulative\_percentage\_series }\OperatorTok{=}\NormalTok{ chocolate\_data[}\StringTok{"cumulative\_percentage"}\NormalTok{]}
\NormalTok{threshold\_index }\OperatorTok{=}\NormalTok{ cumulative\_percentage\_series[cumulative\_percentage\_series }\OperatorTok{\textless{}=}\NormalTok{ threshold].index[}\OperatorTok{{-}}\DecValTok{1}\NormalTok{]}


\NormalTok{y\_start }\OperatorTok{=}\NormalTok{ bars[}\DecValTok{1}\NormalTok{].get\_y()}


\NormalTok{y\_end }\OperatorTok{=} \VariableTok{None}
\ControlFlowTok{for}\NormalTok{ bar, speaker }\KeywordTok{in} \BuiltInTok{zip}\NormalTok{(bars, chocolate\_data[}\StringTok{"speaker"}\NormalTok{]):}
    \ControlFlowTok{if}\NormalTok{ speaker }\OperatorTok{==}\NormalTok{ threshold\_index:}
\NormalTok{        y\_end }\OperatorTok{=}\NormalTok{ bar.get\_y() }\OperatorTok{+}\NormalTok{ bar.get\_height()}
        \ControlFlowTok{break}

\ControlFlowTok{if}\NormalTok{ y\_end }\KeywordTok{is} \VariableTok{None}\NormalTok{:}
    \BuiltInTok{print}\NormalTok{(}\SpecialStringTok{f"No se encontró la barra para el índice: }\SpecialCharTok{\{}\NormalTok{threshold\_index}\SpecialCharTok{\}}\SpecialStringTok{"}\NormalTok{)}
\NormalTok{    y\_end }\OperatorTok{=}\NormalTok{ bars[}\OperatorTok{{-}}\DecValTok{1}\NormalTok{].get\_y() }\OperatorTok{+}\NormalTok{ bars[}\OperatorTok{{-}}\DecValTok{1}\NormalTok{].get\_height()}

\NormalTok{x\_position }\OperatorTok{=}\NormalTok{ chocolate\_data[}\StringTok{"count"}\NormalTok{].}\BuiltInTok{max}\NormalTok{() }\OperatorTok{+} \DecValTok{5}  


\NormalTok{plt.plot([x\_position, x\_position], [y\_start, y\_end], color}\OperatorTok{=}\StringTok{"red"}\NormalTok{, linewidth}\OperatorTok{=}\DecValTok{2}\NormalTok{)  }
\NormalTok{plt.plot([x\_position, x\_position }\OperatorTok{{-}} \DecValTok{10}\NormalTok{], [y\_start, y\_start], color}\OperatorTok{=}\StringTok{"red"}\NormalTok{, linewidth}\OperatorTok{=}\DecValTok{2}\NormalTok{)  }
\NormalTok{plt.plot([x\_position, x\_position }\OperatorTok{{-}} \DecValTok{10}\NormalTok{], [y\_end, y\_end], color}\OperatorTok{=}\StringTok{"red"}\NormalTok{, linewidth}\OperatorTok{=}\DecValTok{2}\NormalTok{)  }


\NormalTok{plt.text(}
\NormalTok{    x\_position }\OperatorTok{+} \DecValTok{5}\NormalTok{, (y\_start }\OperatorTok{+}\NormalTok{ y\_end) }\OperatorTok{/} \DecValTok{2}\NormalTok{, }
    \StringTok{"65.4\% acumulado"}\NormalTok{, fontsize}\OperatorTok{=}\DecValTok{12}\NormalTok{, color}\OperatorTok{=}\StringTok{"red"}\NormalTok{, va}\OperatorTok{=}\StringTok{"center"}\NormalTok{, ha}\OperatorTok{=}\StringTok{"left"}
\NormalTok{)}


\NormalTok{plt.title(}\StringTok{"Cantidad de Discursos por Candidato (Top 5 vs Otros)"}\NormalTok{, fontsize}\OperatorTok{=}\DecValTok{16}\NormalTok{)}
\NormalTok{plt.xlabel(}\StringTok{"Cantidad de Discursos"}\NormalTok{, fontsize}\OperatorTok{=}\DecValTok{14}\NormalTok{)}
\NormalTok{plt.ylabel(}\StringTok{"Candidato"}\NormalTok{, fontsize}\OperatorTok{=}\DecValTok{14}\NormalTok{)}
\NormalTok{plt.grid(axis}\OperatorTok{=}\StringTok{"x"}\NormalTok{, linestyle}\OperatorTok{=}\StringTok{"{-}{-}"}\NormalTok{, alpha}\OperatorTok{=}\FloatTok{0.7}\NormalTok{)}
\NormalTok{plt.tight\_layout()}
\NormalTok{plt.show()}
\end{Highlighting}
\end{Shaded}

\begin{verbatim}
C:\Users\campo\AppData\Local\Temp\ipykernel_5324\4045876432.py:6: SettingWithCopyWarning: 
A value is trying to be set on a copy of a slice from a DataFrame.
Try using .loc[row_indexer,col_indexer] = value instead

See the caveats in the documentation: https://pandas.pydata.org/pandas-docs/stable/user_guide/indexing.html#returning-a-view-versus-a-copy
  top_speeches["month"] = top_speeches["date"].dt.to_period("M")
C:\Users\campo\AppData\Local\Temp\ipykernel_5324\4045876432.py:13: SettingWithCopyWarning: 
A value is trying to be set on a copy of a slice from a DataFrame.
Try using .loc[row_indexer,col_indexer] = value instead

See the caveats in the documentation: https://pandas.pydata.org/pandas-docs/stable/user_guide/indexing.html#returning-a-view-versus-a-copy
  top_speeches["date"] = pd.to_datetime(top_speeches["date"], errors="coerce")
C:\Users\campo\AppData\Local\Temp\ipykernel_5324\4045876432.py:14: SettingWithCopyWarning: 
A value is trying to be set on a copy of a slice from a DataFrame.
Try using .loc[row_indexer,col_indexer] = value instead

See the caveats in the documentation: https://pandas.pydata.org/pandas-docs/stable/user_guide/indexing.html#returning-a-view-versus-a-copy
  top_speeches["month"] = top_speeches["date"].dt.to_period("M")
\end{verbatim}

\begin{verbatim}
No se encontró la barra para el índice: 0
\end{verbatim}

\pandocbounded{\includegraphics[keepaspectratio]{tarea1_campon_loprete_files/figure-pdf/cell-10-output-3.png}}

Los 5 candidatos con mas discursos de campaña representan el 65\% de
todos los discursos registrados y se componen por los oradores Kamala
Harris, Bernie Sanders, Mike Pence, Donald Trump y Joe Biden, dejando de
lado 53 oradores que componen el 35\% restante de discursos registrados.
En base a los discursos de estos 5 candidatos realizaremos el análisis
de aquí en más.

\section{B) Gráfica para visualizar discursos de los candidatos por
año}\label{b-gruxe1fica-para-visualizar-discursos-de-los-candidatos-por-auxf1o}

Si consideramos los discursos a lo largo del tiempo podemos identificar
momentos clave de campaña de los mismos, para ello construimos la figura
**\_\_** que constituye un gráfico de líneas y puntos a lo largo de un
eje temporal.

\begin{Shaded}
\begin{Highlighting}[]
\NormalTok{plt.figure(figsize}\OperatorTok{=}\NormalTok{(}\DecValTok{14}\NormalTok{, }\DecValTok{8}\NormalTok{))}
\NormalTok{sns.lineplot(data}\OperatorTok{=}\NormalTok{full\_merged, x}\OperatorTok{=}\StringTok{"month\_plot"}\NormalTok{, y}\OperatorTok{=}\StringTok{"count"}\NormalTok{, hue}\OperatorTok{=}\StringTok{"speaker"}\NormalTok{, marker}\OperatorTok{=}\StringTok{"o"}\NormalTok{, palette}\OperatorTok{=}\NormalTok{custom\_palette, linewidth}\OperatorTok{=}\FloatTok{2.5}\NormalTok{, markersize}\OperatorTok{=}\DecValTok{8}\NormalTok{, alpha}\OperatorTok{=}\FloatTok{0.65}\NormalTok{)}
\NormalTok{plt.title(}\StringTok{"Cantidad de Discursos por Candidato a lo Largo del Tiempo"}\NormalTok{, fontsize}\OperatorTok{=}\DecValTok{16}\NormalTok{)}
\NormalTok{plt.xlabel(}\StringTok{"Fecha"}\NormalTok{, fontsize}\OperatorTok{=}\DecValTok{14}\NormalTok{)}
\NormalTok{plt.ylabel(}\StringTok{"Cantidad de Discursos"}\NormalTok{, fontsize}\OperatorTok{=}\DecValTok{14}\NormalTok{)}
\NormalTok{plt.legend(title}\OperatorTok{=}\StringTok{"Candidato"}\NormalTok{, fontsize}\OperatorTok{=}\DecValTok{12}\NormalTok{)}
\NormalTok{plt.xticks(rotation}\OperatorTok{=}\DecValTok{0}\NormalTok{, fontsize}\OperatorTok{=}\DecValTok{12}\NormalTok{)}
\NormalTok{plt.yticks(}\BuiltInTok{range}\NormalTok{(}\DecValTok{0}\NormalTok{, np.}\BuiltInTok{max}\NormalTok{(full\_merged[}\StringTok{"count"}\NormalTok{]) }\OperatorTok{+} \DecValTok{1}\NormalTok{, }\DecValTok{1}\NormalTok{), fontsize}\OperatorTok{=}\DecValTok{12}\NormalTok{)}
\NormalTok{plt.grid(axis}\OperatorTok{=}\StringTok{"y"}\NormalTok{, linestyle}\OperatorTok{=}\StringTok{"{-}{-}"}\NormalTok{, alpha}\OperatorTok{=}\FloatTok{0.7}\NormalTok{)}
\NormalTok{eventos }\OperatorTok{=}\NormalTok{ [}
\NormalTok{    \{}\StringTok{"fecha"}\NormalTok{: }\StringTok{"2020{-}04"}\NormalTok{, }\StringTok{"texto"}\NormalTok{: }\StringTok{"Retiro de Bernie Sanders"}\NormalTok{, }\StringTok{"altura"}\NormalTok{: }\FloatTok{0.2}\NormalTok{\},}
\NormalTok{    \{}\StringTok{"fecha"}\NormalTok{: }\StringTok{"2020{-}07"}\NormalTok{, }\StringTok{"texto"}\NormalTok{: }\StringTok{"Inicio campaña Kamala Harris"}\NormalTok{, }\StringTok{"altura"}\NormalTok{: }\FloatTok{0.5}\NormalTok{\},}
\NormalTok{    \{}\StringTok{"fecha"}\NormalTok{: }\StringTok{"2020{-}08"}\NormalTok{, }\StringTok{"texto"}\NormalTok{: }\StringTok{"Formulas: }\CharTok{\textbackslash{}n}\StringTok{Biden \& Harris}\CharTok{\textbackslash{}n}\StringTok{Trump \& Mike Pence"}\NormalTok{, }\StringTok{"altura"}\NormalTok{: }\FloatTok{0.7}\NormalTok{\},}
\NormalTok{    \{}\StringTok{"fecha"}\NormalTok{: }\StringTok{"2020{-}09"}\NormalTok{, }\StringTok{"texto"}\NormalTok{: }\StringTok{"Debate presidencial en TV"}\NormalTok{, }\StringTok{"altura"}\NormalTok{: }\FloatTok{0.9}\NormalTok{\},}
\NormalTok{]}

\ControlFlowTok{for}\NormalTok{ evento }\KeywordTok{in}\NormalTok{ eventos:}
\NormalTok{    fecha }\OperatorTok{=}\NormalTok{ pd.Timestamp(evento[}\StringTok{"fecha"}\NormalTok{])}
\NormalTok{    plt.axvline(fecha, color}\OperatorTok{=}\StringTok{"black"}\NormalTok{, linestyle}\OperatorTok{=}\StringTok{"{-}{-}"}\NormalTok{, alpha}\OperatorTok{=}\FloatTok{0.6}\NormalTok{, linewidth}\OperatorTok{=}\DecValTok{2}\NormalTok{)}
\NormalTok{    plt.text(fecha, plt.ylim()[}\DecValTok{1}\NormalTok{]}\OperatorTok{*}\NormalTok{evento[}\StringTok{"altura"}\NormalTok{], evento[}\StringTok{"texto"}\NormalTok{],}
\NormalTok{             rotation}\OperatorTok{=}\DecValTok{0}\NormalTok{, fontsize}\OperatorTok{=}\DecValTok{11}\NormalTok{, color}\OperatorTok{=}\StringTok{"black"}\NormalTok{,}
\NormalTok{             verticalalignment}\OperatorTok{=}\StringTok{"top"}\NormalTok{, horizontalalignment}\OperatorTok{=}\StringTok{"left"}\NormalTok{,}
\NormalTok{             bbox}\OperatorTok{=}\BuiltInTok{dict}\NormalTok{(boxstyle}\OperatorTok{=}\StringTok{"round,pad=0.3"}\NormalTok{, facecolor}\OperatorTok{=}\StringTok{"white"}\NormalTok{, edgecolor}\OperatorTok{=}\StringTok{"gray"}\NormalTok{, alpha}\OperatorTok{=}\FloatTok{0.8}\NormalTok{))}

\NormalTok{plt.tight\_layout()}
\NormalTok{plt.show()}
\end{Highlighting}
\end{Shaded}

\pandocbounded{\includegraphics[keepaspectratio]{tarea1_campon_loprete_files/figure-pdf/cell-11-output-1.png}}

Como se puede ver en la figura anterior, la campaña de los candidatos
arranca en Enero 2020, con pocos discursos mensuales, no más de 5. En
Abril se observa como la cantidad de discursos de Bernie bajan a 0, lo
cual coincide con su anuncio de retirada de la campaña y su apoyo a Joe
Biden por parte del partido Demócrata de EE.UU. por eso se explica que
se visualiza un único discurso de este candidato en los últimos 3 meses
de campaña. Por otro lado, en Julio se observa el inicio de camapaña de
Kamala Harris y posteriormente en Agosto su anuncio de formula
presidencial junto a Joe Biden como presidente. En Septiembre se alcanza
el pico de discursos de los 2 candidatos principales a preseidente de
los EE.UU. producto del debate presidencial que es trasmitido en todo el
país y la llegada de las elecciones en el mes siguiente.

La figura **\_\_** muestra los discursos mensuales de las formulas
presidenciales, este gráfico muestra de forma agregada los discursos de
Joe Biden \& Kamala Harris y Donald Trump \& Mike Pence; obsérvese como
en Julio comienza a incrementarse la cantidad de discursos por fórmula.

\begin{Shaded}
\begin{Highlighting}[]
\CommentTok{\# creamos las fórmulas presidenciales}
\NormalTok{formula\_map }\OperatorTok{=}\NormalTok{ \{}
    \StringTok{"Joe Biden"}\NormalTok{: }\StringTok{"Joe Biden \& Kamala Harris"}\NormalTok{,}
    \StringTok{"Kamala Harris"}\NormalTok{: }\StringTok{"Joe Biden \& Kamala Harris"}\NormalTok{,}
    \StringTok{"Donald Trump"}\NormalTok{: }\StringTok{"Donald Trump \& Mike Pence"}\NormalTok{,}
    \StringTok{"Mike Pence"}\NormalTok{: }\StringTok{"Donald Trump \& Mike Pence"}
\NormalTok{\}}

\NormalTok{full\_merged[}\StringTok{"formula"}\NormalTok{] }\OperatorTok{=}\NormalTok{ full\_merged[}\StringTok{"speaker"}\NormalTok{].}\BuiltInTok{map}\NormalTok{(formula\_map)}
\CommentTok{\# filtramos a sanders}
\NormalTok{formula\_df }\OperatorTok{=}\NormalTok{ full\_merged[full\_merged[}\StringTok{"formula"}\NormalTok{].notnull()]}

\NormalTok{formula\_grouped }\OperatorTok{=}\NormalTok{ formula\_df.groupby([}\StringTok{"month"}\NormalTok{, }\StringTok{"formula"}\NormalTok{]).agg(\{}
    \StringTok{"count"}\NormalTok{: }\StringTok{"sum"}
\NormalTok{\}).reset\_index()}
\NormalTok{formula\_grouped[}\StringTok{"month\_plot"}\NormalTok{] }\OperatorTok{=}\NormalTok{ formula\_grouped[}\StringTok{"month"}\NormalTok{].dt.to\_timestamp()}

\NormalTok{plt.figure(figsize}\OperatorTok{=}\NormalTok{(}\DecValTok{14}\NormalTok{, }\DecValTok{8}\NormalTok{))}
\NormalTok{sns.lineplot(data}\OperatorTok{=}\NormalTok{formula\_grouped, x}\OperatorTok{=}\StringTok{"month\_plot"}\NormalTok{, y}\OperatorTok{=}\StringTok{"count"}\NormalTok{, hue}\OperatorTok{=}\StringTok{"formula"}\NormalTok{, marker}\OperatorTok{=}\StringTok{"o"}\NormalTok{, linewidth}\OperatorTok{=}\FloatTok{3.5}\NormalTok{, markersize}\OperatorTok{=}\DecValTok{10}\NormalTok{)}

\NormalTok{plt.title(}\StringTok{"Cantidad de Discursos por Fórmula Presidencial"}\NormalTok{, fontsize}\OperatorTok{=}\DecValTok{16}\NormalTok{)}
\NormalTok{plt.xlabel(}\StringTok{"Fecha"}\NormalTok{, fontsize}\OperatorTok{=}\DecValTok{14}\NormalTok{)}
\NormalTok{plt.ylabel(}\StringTok{"Cantidad de Discursos"}\NormalTok{, fontsize}\OperatorTok{=}\DecValTok{14}\NormalTok{)}
\NormalTok{plt.legend(title}\OperatorTok{=}\StringTok{"Fórmula Presidencial"}\NormalTok{, fontsize}\OperatorTok{=}\DecValTok{12}\NormalTok{)}
\NormalTok{plt.xticks(rotation}\OperatorTok{=}\DecValTok{0}\NormalTok{, fontsize}\OperatorTok{=}\DecValTok{12}\NormalTok{)}
\NormalTok{plt.yticks(}\BuiltInTok{range}\NormalTok{(}\DecValTok{0}\NormalTok{, formula\_grouped[}\StringTok{"count"}\NormalTok{].}\BuiltInTok{max}\NormalTok{() }\OperatorTok{+} \DecValTok{1}\NormalTok{), fontsize}\OperatorTok{=}\DecValTok{12}\NormalTok{)}
\NormalTok{plt.grid(axis}\OperatorTok{=}\StringTok{"y"}\NormalTok{, linestyle}\OperatorTok{=}\StringTok{"{-}{-}"}\NormalTok{, alpha}\OperatorTok{=}\FloatTok{0.7}\NormalTok{)}

\NormalTok{eventos }\OperatorTok{=}\NormalTok{ [}
\NormalTok{    \{}\StringTok{"fecha"}\NormalTok{: }\StringTok{"2020{-}04"}\NormalTok{, }\StringTok{"texto"}\NormalTok{: }\StringTok{"Retiro de Bernie Sanders"}\NormalTok{, }\StringTok{"altura"}\NormalTok{: }\FloatTok{0.2}\NormalTok{\},}
\NormalTok{    \{}\StringTok{"fecha"}\NormalTok{: }\StringTok{"2020{-}07"}\NormalTok{, }\StringTok{"texto"}\NormalTok{: }\StringTok{"Inicio campaña Kamala Harris"}\NormalTok{, }\StringTok{"altura"}\NormalTok{: }\FloatTok{0.5}\NormalTok{\},}
\NormalTok{    \{}\StringTok{"fecha"}\NormalTok{: }\StringTok{"2020{-}08"}\NormalTok{, }\StringTok{"texto"}\NormalTok{: }\StringTok{"Formulas: }\CharTok{\textbackslash{}n}\StringTok{Biden \& Harris}\CharTok{\textbackslash{}n}\StringTok{Trump \& Mike Pence"}\NormalTok{, }\StringTok{"altura"}\NormalTok{: }\FloatTok{0.7}\NormalTok{\},}
\NormalTok{    \{}\StringTok{"fecha"}\NormalTok{: }\StringTok{"2020{-}09"}\NormalTok{, }\StringTok{"texto"}\NormalTok{: }\StringTok{"Debate presidencial en TV"}\NormalTok{, }\StringTok{"altura"}\NormalTok{: }\FloatTok{0.9}\NormalTok{\},}
\NormalTok{]}

\ControlFlowTok{for}\NormalTok{ evento }\KeywordTok{in}\NormalTok{ eventos:}
\NormalTok{    fecha }\OperatorTok{=}\NormalTok{ pd.Timestamp(evento[}\StringTok{"fecha"}\NormalTok{])}
\NormalTok{    plt.axvline(fecha, color}\OperatorTok{=}\StringTok{"black"}\NormalTok{, linestyle}\OperatorTok{=}\StringTok{"{-}{-}"}\NormalTok{, alpha}\OperatorTok{=}\FloatTok{0.6}\NormalTok{, linewidth}\OperatorTok{=}\DecValTok{2}\NormalTok{)}
\NormalTok{    plt.text(fecha, plt.ylim()[}\DecValTok{1}\NormalTok{]}\OperatorTok{*}\NormalTok{evento[}\StringTok{"altura"}\NormalTok{], evento[}\StringTok{"texto"}\NormalTok{],}
\NormalTok{             rotation}\OperatorTok{=}\DecValTok{0}\NormalTok{, fontsize}\OperatorTok{=}\DecValTok{11}\NormalTok{, color}\OperatorTok{=}\StringTok{"black"}\NormalTok{,}
\NormalTok{             verticalalignment}\OperatorTok{=}\StringTok{"top"}\NormalTok{, horizontalalignment}\OperatorTok{=}\StringTok{"left"}\NormalTok{,}
\NormalTok{             bbox}\OperatorTok{=}\BuiltInTok{dict}\NormalTok{(boxstyle}\OperatorTok{=}\StringTok{"round,pad=0.3"}\NormalTok{, facecolor}\OperatorTok{=}\StringTok{"white"}\NormalTok{, edgecolor}\OperatorTok{=}\StringTok{"gray"}\NormalTok{, alpha}\OperatorTok{=}\FloatTok{0.8}\NormalTok{))}

\NormalTok{plt.tight\_layout()}
\NormalTok{plt.show()}
\end{Highlighting}
\end{Shaded}

\pandocbounded{\includegraphics[keepaspectratio]{tarea1_campon_loprete_files/figure-pdf/cell-12-output-1.png}}

\begin{Shaded}
\begin{Highlighting}[]
\NormalTok{top\_speeches[}\StringTok{"date"}\NormalTok{] }\OperatorTok{=}\NormalTok{ pd.to\_datetime(top\_speeches[}\StringTok{"date"}\NormalTok{], errors}\OperatorTok{=}\StringTok{"coerce"}\NormalTok{)}

\CommentTok{\# asignaciones de bimestres}
\KeywordTok{def}\NormalTok{ assign\_bimester(date):}
\NormalTok{    month }\OperatorTok{=}\NormalTok{ date.month}
    \ControlFlowTok{if}\NormalTok{ month }\KeywordTok{in}\NormalTok{ [}\DecValTok{1}\NormalTok{, }\DecValTok{2}\NormalTok{]:}
        \ControlFlowTok{return} \StringTok{"Ene{-}Feb"}
    \ControlFlowTok{elif}\NormalTok{ month }\KeywordTok{in}\NormalTok{ [}\DecValTok{3}\NormalTok{, }\DecValTok{4}\NormalTok{]:}
        \ControlFlowTok{return} \StringTok{"Mar{-}Abr"}
    \ControlFlowTok{elif}\NormalTok{ month }\KeywordTok{in}\NormalTok{ [}\DecValTok{5}\NormalTok{, }\DecValTok{6}\NormalTok{]:}
        \ControlFlowTok{return} \StringTok{"May{-}Jun"}
    \ControlFlowTok{elif}\NormalTok{ month }\KeywordTok{in}\NormalTok{ [}\DecValTok{7}\NormalTok{, }\DecValTok{8}\NormalTok{]:}
        \ControlFlowTok{return} \StringTok{"Jul{-}Ago"}
    \ControlFlowTok{elif}\NormalTok{ month }\KeywordTok{in}\NormalTok{ [}\DecValTok{9}\NormalTok{, }\DecValTok{10}\NormalTok{]:}
        \ControlFlowTok{return} \StringTok{"Sep{-}Oct"}
    \ControlFlowTok{else}\NormalTok{:}
        \ControlFlowTok{return} \StringTok{"Nov{-}Dic"}  \CommentTok{\# Opcional, si hay datos en nov{-}dic}


\NormalTok{top\_speeches[}\StringTok{"bimester"}\NormalTok{] }\OperatorTok{=}\NormalTok{ top\_speeches[}\StringTok{"date"}\NormalTok{].}\BuiltInTok{apply}\NormalTok{(assign\_bimester)}

\NormalTok{bimester\_data }\OperatorTok{=}\NormalTok{ top\_speeches.groupby([}\StringTok{"bimester"}\NormalTok{, }\StringTok{"speaker"}\NormalTok{]).size().reset\_index(name}\OperatorTok{=}\StringTok{"count"}\NormalTok{)}

\CommentTok{\# ordenamos los bimestres}
\NormalTok{ordered\_bimesters }\OperatorTok{=}\NormalTok{ [}\StringTok{"Ene{-}Feb"}\NormalTok{, }\StringTok{"Mar{-}Abr"}\NormalTok{, }\StringTok{"May{-}Jun"}\NormalTok{, }\StringTok{"Jul{-}Ago"}\NormalTok{, }\StringTok{"Sep{-}Oct"}\NormalTok{]}
\NormalTok{bimester\_data[}\StringTok{"bimester"}\NormalTok{] }\OperatorTok{=}\NormalTok{ pd.Categorical(bimester\_data[}\StringTok{"bimester"}\NormalTok{], categories}\OperatorTok{=}\NormalTok{ordered\_bimesters, ordered}\OperatorTok{=}\VariableTok{True}\NormalTok{)}
\NormalTok{bimester\_data }\OperatorTok{=}\NormalTok{ bimester\_data.sort\_values([}\StringTok{"bimester"}\NormalTok{, }\StringTok{"count"}\NormalTok{], ascending}\OperatorTok{=}\NormalTok{[}\VariableTok{True}\NormalTok{, }\VariableTok{False}\NormalTok{])}

\NormalTok{fig, axes }\OperatorTok{=}\NormalTok{ plt.subplots(nrows}\OperatorTok{=}\DecValTok{5}\NormalTok{, ncols}\OperatorTok{=}\DecValTok{1}\NormalTok{, figsize}\OperatorTok{=}\NormalTok{(}\DecValTok{10}\NormalTok{, }\DecValTok{18}\NormalTok{), sharex}\OperatorTok{=}\VariableTok{True}\NormalTok{)}

\ControlFlowTok{for}\NormalTok{ i, bim }\KeywordTok{in} \BuiltInTok{enumerate}\NormalTok{(ordered\_bimesters):}
\NormalTok{    ax }\OperatorTok{=}\NormalTok{ axes[i]}
\NormalTok{    data }\OperatorTok{=}\NormalTok{ bimester\_data[bimester\_data[}\StringTok{"bimester"}\NormalTok{] }\OperatorTok{==}\NormalTok{ bim].sort\_values(by}\OperatorTok{=}\StringTok{"count"}\NormalTok{, ascending}\OperatorTok{=}\VariableTok{True}\NormalTok{)}
    
\NormalTok{    ax.barh(data[}\StringTok{"speaker"}\NormalTok{], data[}\StringTok{"count"}\NormalTok{], color}\OperatorTok{=}\StringTok{"skyblue"}\NormalTok{, alpha}\OperatorTok{=}\FloatTok{0.7}\NormalTok{)}
\NormalTok{    ax.set\_title(}\SpecialStringTok{f"Bimestre: }\SpecialCharTok{\{}\NormalTok{bim}\SpecialCharTok{\}}\SpecialStringTok{"}\NormalTok{)}
\NormalTok{    ax.set\_xlabel(}\StringTok{"Cantidad de Discursos"}\NormalTok{)}
\NormalTok{    ax.set\_ylabel(}\StringTok{"Candidato"}\NormalTok{)}
\NormalTok{    ax.set\_xlim(left}\OperatorTok{=}\DecValTok{0}\NormalTok{)}
\NormalTok{    ax.set\_xticks(np.arange(}\DecValTok{0}\NormalTok{, data[}\StringTok{"count"}\NormalTok{].}\BuiltInTok{max}\NormalTok{() }\OperatorTok{+} \DecValTok{1}\NormalTok{))}

\NormalTok{plt.tight\_layout()}
\CommentTok{\#plt.show()}
\NormalTok{plt.close()}
\end{Highlighting}
\end{Shaded}

\begin{verbatim}
C:\Users\campo\AppData\Local\Temp\ipykernel_5324\1495649151.py:1: SettingWithCopyWarning: 
A value is trying to be set on a copy of a slice from a DataFrame.
Try using .loc[row_indexer,col_indexer] = value instead

See the caveats in the documentation: https://pandas.pydata.org/pandas-docs/stable/user_guide/indexing.html#returning-a-view-versus-a-copy
  top_speeches["date"] = pd.to_datetime(top_speeches["date"], errors="coerce")
C:\Users\campo\AppData\Local\Temp\ipykernel_5324\1495649151.py:20: SettingWithCopyWarning: 
A value is trying to be set on a copy of a slice from a DataFrame.
Try using .loc[row_indexer,col_indexer] = value instead

See the caveats in the documentation: https://pandas.pydata.org/pandas-docs/stable/user_guide/indexing.html#returning-a-view-versus-a-copy
  top_speeches["bimester"] = top_speeches["date"].apply(assign_bimester)
\end{verbatim}

\section{Parte C) Limpieza de texto}\label{parte-c-limpieza-de-texto}

\begin{Shaded}
\begin{Highlighting}[]
\KeywordTok{def}\NormalTok{ clean\_text(df, column\_name):}
\NormalTok{    result }\OperatorTok{=}\NormalTok{ df[column\_name]}

    \CommentTok{\# quitar todo hasta el primer \textbackslash{}n}
\NormalTok{    result }\OperatorTok{=}\NormalTok{ result.}\BuiltInTok{str}\NormalTok{.replace(}\VerbatimStringTok{r"}\DecValTok{\^{}}\PreprocessorTok{[\^{}}\CharTok{\textbackslash{}n}\PreprocessorTok{]}\OperatorTok{*}\CharTok{\textbackslash{}n}\VerbatimStringTok{"}\NormalTok{, }\StringTok{""}\NormalTok{, regex}\OperatorTok{=}\VariableTok{True}\NormalTok{)}

    \CommentTok{\# convertir a minúsculas}
\NormalTok{    result }\OperatorTok{=}\NormalTok{ result.}\BuiltInTok{str}\NormalTok{.lower()}

    \CommentTok{\# eliminar signos de puntuación definidos}
\NormalTok{    result }\OperatorTok{=}\NormalTok{ result.}\BuiltInTok{str}\NormalTok{.replace(}\VerbatimStringTok{r"}\PreprocessorTok{[\{\}}\CharTok{\textbackslash{}n}\PreprocessorTok{]}\VerbatimStringTok{"}\NormalTok{.}\BuiltInTok{format}\NormalTok{(re.escape(}\StringTok{"[],:?.!;}\CharTok{\textbackslash{}"}\StringTok{\textquotesingle{}"}\NormalTok{)), }\StringTok{" "}\NormalTok{, regex}\OperatorTok{=}\VariableTok{True}\NormalTok{)}

    \CommentTok{\# eliminar espacios extra}
\NormalTok{    result }\OperatorTok{=}\NormalTok{ result.}\BuiltInTok{str}\NormalTok{.replace(}\VerbatimStringTok{r"}\DecValTok{\textbackslash{}s}\OperatorTok{+}\VerbatimStringTok{"}\NormalTok{, }\StringTok{" "}\NormalTok{, regex}\OperatorTok{=}\VariableTok{True}\NormalTok{).}\BuiltInTok{str}\NormalTok{.strip()}

    \ControlFlowTok{return}\NormalTok{ result}

\CommentTok{\# creamos columna clentext}
\NormalTok{top\_speeches[}\StringTok{"CleanText"}\NormalTok{] }\OperatorTok{=}\NormalTok{ clean\_text(top\_speeches, }\StringTok{"text"}\NormalTok{)}

\CommentTok{\# Convierte párrafos en listas "palabra1 palabra2 palabra3" {-}\textgreater{} ["palabra1", "palabra2", "palabra3"]}
\NormalTok{top\_speeches[}\StringTok{"WordList"}\NormalTok{] }\OperatorTok{=}\NormalTok{ top\_speeches[}\StringTok{"CleanText"}\NormalTok{].}\BuiltInTok{str}\NormalTok{.split()}

\CommentTok{\# Veamos la nueva columna creada: notar que a la derecha tenemos una lista: [palabra1, palabra2, palabra3]}
\CommentTok{\#top\_speeches[["CleanText", "WordList"]]}
\end{Highlighting}
\end{Shaded}

\begin{verbatim}
C:\Users\campo\AppData\Local\Temp\ipykernel_5324\23623688.py:19: SettingWithCopyWarning: 
A value is trying to be set on a copy of a slice from a DataFrame.
Try using .loc[row_indexer,col_indexer] = value instead

See the caveats in the documentation: https://pandas.pydata.org/pandas-docs/stable/user_guide/indexing.html#returning-a-view-versus-a-copy
  top_speeches["CleanText"] = clean_text(top_speeches, "text")
C:\Users\campo\AppData\Local\Temp\ipykernel_5324\23623688.py:22: SettingWithCopyWarning: 
A value is trying to be set on a copy of a slice from a DataFrame.
Try using .loc[row_indexer,col_indexer] = value instead

See the caveats in the documentation: https://pandas.pydata.org/pandas-docs/stable/user_guide/indexing.html#returning-a-view-versus-a-copy
  top_speeches["WordList"] = top_speeches["CleanText"].str.split()
\end{verbatim}

Para normalizar el texto realizamos unas pequeñas modificaciones en la
función de \texttt{clean_text()}, agregamos una sentencia que elimina
los signos de puntuacion definidos y los reemplaza por espacios.
Específicamente, eliminamos los caracteres: {[} {]} , : ? . ! ; '' ',
además de saltos de línea (\textbackslash n), usando una expresión
regular. Este proceso nos permite limpiar el texto antes de realizar
análisis de texto.

\section{Parte 2)}\label{parte-2}

\subsection{Parte A)}\label{parte-a}

\begin{Shaded}
\begin{Highlighting}[]
\CommentTok{\# stopwords es un diccionario de palabras que no aportan significado}
\NormalTok{nltk.download(}\StringTok{\textquotesingle{}stopwords\textquotesingle{}}\NormalTok{)}
\NormalTok{stop\_words }\OperatorTok{=} \BuiltInTok{set}\NormalTok{(stopwords.words(}\StringTok{\textquotesingle{}english\textquotesingle{}}\NormalTok{))}

\ControlFlowTok{if} \StringTok{"WordList"} \KeywordTok{not} \KeywordTok{in}\NormalTok{ top\_speeches.columns:}
\NormalTok{    top\_speeches[}\StringTok{"WordList"}\NormalTok{] }\OperatorTok{=}\NormalTok{ top\_speeches[}\StringTok{"CleanText"}\NormalTok{].}\BuiltInTok{str}\NormalTok{.split()}

\NormalTok{custom\_stopwords }\OperatorTok{=}\NormalTok{ \{}\StringTok{"it’s"}\NormalTok{\}}

\NormalTok{top\_speeches[}\StringTok{"FilteredWordList"}\NormalTok{] }\OperatorTok{=}\NormalTok{ top\_speeches[}\StringTok{"WordList"}\NormalTok{].}\BuiltInTok{apply}\NormalTok{(}
    \KeywordTok{lambda}\NormalTok{ words: [}
\NormalTok{        word }\ControlFlowTok{for}\NormalTok{ word }\KeywordTok{in}\NormalTok{ words }
        \ControlFlowTok{if}\NormalTok{ word }\KeywordTok{not} \KeywordTok{in}\NormalTok{ stop\_words }\KeywordTok{and}\NormalTok{ word }\KeywordTok{not} \KeywordTok{in}\NormalTok{ custom\_stopwords }\KeywordTok{and} \BuiltInTok{len}\NormalTok{(word) }\OperatorTok{\textgreater{}} \DecValTok{2}
\NormalTok{    ]}
\NormalTok{)}

\CommentTok{\# contar palabras por candidato}
\NormalTok{word\_counts }\OperatorTok{=}\NormalTok{ \{\}}
\ControlFlowTok{for}\NormalTok{ speaker }\KeywordTok{in}\NormalTok{ top\_speeches[}\StringTok{"speaker"}\NormalTok{].unique():}
\NormalTok{    all\_words }\OperatorTok{=}\NormalTok{ top\_speeches[top\_speeches[}\StringTok{"speaker"}\NormalTok{] }\OperatorTok{==}\NormalTok{ speaker][}\StringTok{"FilteredWordList"}\NormalTok{].explode()}
\NormalTok{    word\_counts[speaker] }\OperatorTok{=}\NormalTok{ Counter(all\_words)}

\CommentTok{\# crear un df con las palabras más frecuentes}
\NormalTok{top\_words }\OperatorTok{=}\NormalTok{ pd.DataFrame([}
\NormalTok{    \{}\StringTok{"speaker"}\NormalTok{: speaker, }\StringTok{"word"}\NormalTok{: word, }\StringTok{"count"}\NormalTok{: count\}}
    \ControlFlowTok{for}\NormalTok{ speaker, counts }\KeywordTok{in}\NormalTok{ word\_counts.items()}
    \ControlFlowTok{for}\NormalTok{ word, count }\KeywordTok{in}\NormalTok{ counts.most\_common(}\DecValTok{10}\NormalTok{)}
\NormalTok{])}
\CommentTok{\# reordenamos por frecuencia}
\NormalTok{top\_words }\OperatorTok{=}\NormalTok{ top\_words.sort\_values(by}\OperatorTok{=}\NormalTok{[}\StringTok{"count"}\NormalTok{], ascending}\OperatorTok{=}\NormalTok{[}\VariableTok{False}\NormalTok{])}

\NormalTok{plt.figure(figsize}\OperatorTok{=}\NormalTok{(}\DecValTok{12}\NormalTok{, }\DecValTok{8}\NormalTok{))}
\NormalTok{sns.barplot(data}\OperatorTok{=}\NormalTok{top\_words, x}\OperatorTok{=}\StringTok{"count"}\NormalTok{, y}\OperatorTok{=}\StringTok{"word"}\NormalTok{, hue}\OperatorTok{=}\StringTok{"speaker"}\NormalTok{, dodge}\OperatorTok{=}\VariableTok{False}\NormalTok{, palette}\OperatorTok{=}\NormalTok{custom\_palette)}
\NormalTok{plt.title(}\StringTok{"Palabras Más Frecuentes por Candidato"}\NormalTok{)}
\NormalTok{plt.xlabel(}\StringTok{"Frecuencia"}\NormalTok{)}
\NormalTok{plt.ylabel(}\StringTok{"Palabras"}\NormalTok{)}
\NormalTok{plt.legend(title}\OperatorTok{=}\StringTok{"Candidato"}\NormalTok{)}
\NormalTok{plt.tight\_layout()}
\NormalTok{plt.show()}
\end{Highlighting}
\end{Shaded}

\begin{verbatim}
[nltk_data] Downloading package stopwords to
[nltk_data]     C:\Users\campo\AppData\Roaming\nltk_data...
[nltk_data]   Package stopwords is already up-to-date!
C:\Users\campo\AppData\Local\Temp\ipykernel_5324\1985702901.py:10: SettingWithCopyWarning: 
A value is trying to be set on a copy of a slice from a DataFrame.
Try using .loc[row_indexer,col_indexer] = value instead

See the caveats in the documentation: https://pandas.pydata.org/pandas-docs/stable/user_guide/indexing.html#returning-a-view-versus-a-copy
  top_speeches["FilteredWordList"] = top_speeches["WordList"].apply(
\end{verbatim}

\pandocbounded{\includegraphics[keepaspectratio]{tarea1_campon_loprete_files/figure-pdf/cell-15-output-2.png}}

\begin{Shaded}
\begin{Highlighting}[]
\CommentTok{\#ahora construimos grafico de palabras mas frecuentes por formulas presidenciales}
\NormalTok{top\_speeches[}\StringTok{"formula"}\NormalTok{] }\OperatorTok{=}\NormalTok{ top\_speeches[}\StringTok{"speaker"}\NormalTok{].}\BuiltInTok{map}\NormalTok{(formula\_map)}
\NormalTok{formula\_speeches }\OperatorTok{=}\NormalTok{ top\_speeches.dropna(subset}\OperatorTok{=}\NormalTok{[}\StringTok{"formula"}\NormalTok{])}

\NormalTok{formula\_word\_counts }\OperatorTok{=}\NormalTok{ \{\}}
\ControlFlowTok{for}\NormalTok{ formula }\KeywordTok{in}\NormalTok{ formula\_speeches[}\StringTok{"formula"}\NormalTok{].unique():}
\NormalTok{    all\_words }\OperatorTok{=}\NormalTok{ formula\_speeches[formula\_speeches[}\StringTok{"formula"}\NormalTok{] }\OperatorTok{==}\NormalTok{ formula][}\StringTok{"FilteredWordList"}\NormalTok{].explode()}
\NormalTok{    formula\_word\_counts[formula] }\OperatorTok{=}\NormalTok{ Counter(all\_words)}

\NormalTok{top\_formula\_words }\OperatorTok{=}\NormalTok{ pd.DataFrame([}
\NormalTok{    \{}\StringTok{"formula"}\NormalTok{: formula, }\StringTok{"word"}\NormalTok{: word, }\StringTok{"count"}\NormalTok{: count\}}
    \ControlFlowTok{for}\NormalTok{ formula, counts }\KeywordTok{in}\NormalTok{ formula\_word\_counts.items()}
    \ControlFlowTok{for}\NormalTok{ word, count }\KeywordTok{in}\NormalTok{ counts.most\_common(}\DecValTok{10}\NormalTok{)}
\NormalTok{])}

\NormalTok{top\_formula\_words }\OperatorTok{=}\NormalTok{ top\_formula\_words.sort\_values(by}\OperatorTok{=}\StringTok{"count"}\NormalTok{, ascending}\OperatorTok{=}\VariableTok{False}\NormalTok{)}

\NormalTok{plt.figure(figsize}\OperatorTok{=}\NormalTok{(}\DecValTok{12}\NormalTok{, }\DecValTok{8}\NormalTok{))}
\NormalTok{sns.barplot(data}\OperatorTok{=}\NormalTok{top\_formula\_words, x}\OperatorTok{=}\StringTok{"count"}\NormalTok{, y}\OperatorTok{=}\StringTok{"word"}\NormalTok{, hue}\OperatorTok{=}\StringTok{"formula"}\NormalTok{, dodge}\OperatorTok{=}\VariableTok{False}\NormalTok{, palette}\OperatorTok{=}\NormalTok{custom\_palette[:}\DecValTok{2}\NormalTok{])}
\NormalTok{plt.title(}\StringTok{"Palabras Más Frecuentes por Fórmula Presidencial"}\NormalTok{)}
\NormalTok{plt.xlabel(}\StringTok{"Frecuencia"}\NormalTok{)}
\NormalTok{plt.ylabel(}\StringTok{"Palabras"}\NormalTok{)}
\NormalTok{plt.legend(title}\OperatorTok{=}\StringTok{"Fórmula"}\NormalTok{)}
\NormalTok{plt.tight\_layout()}
\NormalTok{plt.show()}
\end{Highlighting}
\end{Shaded}

\begin{verbatim}
C:\Users\campo\AppData\Local\Temp\ipykernel_5324\2785974830.py:2: SettingWithCopyWarning: 
A value is trying to be set on a copy of a slice from a DataFrame.
Try using .loc[row_indexer,col_indexer] = value instead

See the caveats in the documentation: https://pandas.pydata.org/pandas-docs/stable/user_guide/indexing.html#returning-a-view-versus-a-copy
  top_speeches["formula"] = top_speeches["speaker"].map(formula_map)
\end{verbatim}

\pandocbounded{\includegraphics[keepaspectratio]{tarea1_campon_loprete_files/figure-pdf/cell-16-output-2.png}}

\section{TODO: Ordenar todo esto y redactar (Es de la parte
2)}\label{todo-ordenar-todo-esto-y-redactar-es-de-la-parte-2}

\begin{Shaded}
\begin{Highlighting}[]
\NormalTok{candidates }\OperatorTok{=}\NormalTok{ top\_speeches[}\StringTok{"speaker"}\NormalTok{].unique()}
\NormalTok{mentions\_matrix }\OperatorTok{=}\NormalTok{ pd.DataFrame(}\DecValTok{0}\NormalTok{, index}\OperatorTok{=}\NormalTok{candidates, columns}\OperatorTok{=}\NormalTok{candidates)}

\CommentTok{\# Menciones entre candidatos (sin contar a sí mismos) \# FIX: si menciona a mas de uno se cuenta en cada candidato}
\ControlFlowTok{for}\NormalTok{ speaker }\KeywordTok{in}\NormalTok{ candidates:}
    \ControlFlowTok{for}\NormalTok{ mentioned }\KeywordTok{in}\NormalTok{ candidates:}
\NormalTok{        mentions\_matrix.loc[speaker, mentioned] }\OperatorTok{=}\NormalTok{ top\_speeches[}
\NormalTok{            (top\_speeches[}\StringTok{"speaker"}\NormalTok{] }\OperatorTok{==}\NormalTok{ speaker) }\OperatorTok{\&} 
\NormalTok{            (top\_speeches[}\StringTok{"CleanText"}\NormalTok{].}\BuiltInTok{str}\NormalTok{.contains(mentioned, case}\OperatorTok{=}\VariableTok{False}\NormalTok{))}
\NormalTok{        ].shape[}\DecValTok{0}\NormalTok{]}

\CommentTok{\# Visualizar la matriz}
\NormalTok{plt.figure(figsize}\OperatorTok{=}\NormalTok{(}\DecValTok{8}\NormalTok{, }\DecValTok{6}\NormalTok{))}
\NormalTok{sns.heatmap(mentions\_matrix, annot}\OperatorTok{=}\VariableTok{True}\NormalTok{, fmt}\OperatorTok{=}\StringTok{"d"}\NormalTok{, cmap}\OperatorTok{=}\StringTok{"Blues"}\NormalTok{, cbar}\OperatorTok{=}\VariableTok{False}\NormalTok{)}
\NormalTok{plt.title(}\StringTok{"Matriz de Menciones entre Candidatos"}\NormalTok{)}
\NormalTok{plt.xlabel(}\StringTok{"Mencionado"}\NormalTok{)}
\NormalTok{plt.ylabel(}\StringTok{"Candidato"}\NormalTok{)}
\NormalTok{plt.tight\_layout()}
\NormalTok{plt.show()}
\end{Highlighting}
\end{Shaded}

\pandocbounded{\includegraphics[keepaspectratio]{tarea1_campon_loprete_files/figure-pdf/cell-17-output-1.png}}

\begin{Shaded}
\begin{Highlighting}[]
\CommentTok{\# Crear un grafo dirigido}
\NormalTok{G }\OperatorTok{=}\NormalTok{ nx.from\_pandas\_adjacency(mentions\_matrix, create\_using}\OperatorTok{=}\NormalTok{nx.DiGraph)}
\NormalTok{plt.figure(figsize}\OperatorTok{=}\NormalTok{(}\DecValTok{10}\NormalTok{, }\DecValTok{8}\NormalTok{))}
\NormalTok{pos }\OperatorTok{=}\NormalTok{ nx.spring\_layout(G)}
\NormalTok{nx.draw(G, pos, with\_labels}\OperatorTok{=}\VariableTok{True}\NormalTok{, node\_size}\OperatorTok{=}\DecValTok{3000}\NormalTok{, node\_color}\OperatorTok{=}\StringTok{"lightblue"}\NormalTok{, font\_size}\OperatorTok{=}\DecValTok{10}\NormalTok{, font\_weight}\OperatorTok{=}\StringTok{"bold"}\NormalTok{, arrowsize}\OperatorTok{=}\DecValTok{20}\NormalTok{)}
\NormalTok{plt.title(}\StringTok{"Grafo de Menciones entre Candidatos"}\NormalTok{)}
\NormalTok{plt.show()}
\end{Highlighting}
\end{Shaded}

\pandocbounded{\includegraphics[keepaspectratio]{tarea1_campon_loprete_files/figure-pdf/cell-18-output-1.png}}




\end{document}
